\documentclass[12pt]{article}
\usepackage[margin=0.72in]{geometry}
\usepackage{lmodern}
\usepackage{microtype}
\usepackage{multicol}
\usepackage{fancyhdr}
\usepackage{titlesec}
\usepackage{enumitem}
\usepackage{xcolor}
\usepackage[hidelinks]{hyperref}

\setlength{\columnsep}{0.30in}
\setlength{\parindent}{1.05em}
\setlength{\parskip}{0.24em}
\setlength{\headheight}{15pt}
\titleformat{\section}{\large\bfseries\scshape}{}{0pt}{}
\titleformat{\subsection}{\normalsize\bfseries}{}{0pt}{}
\pagestyle{fancy}
\fancyhf{}
\fancyhead[L]{Kant on Enlightenment}
\fancyhead[R]{Reading Handout}
\fancyfoot[C]{\thepage}
\renewcommand{\headrulewidth}{0.4pt}

\newcommand{\focusbox}[1]{%
\vspace{0.4em}\noindent\colorbox{gray!12}{\begin{minipage}{0.96\linewidth}#1\end{minipage}}\vspace{0.5em}}

\begin{document}
\begin{center}
{\LARGE\bfseries Do We Love Shadows?}\par\vspace{0.25em}
{\large\scshape Reading 5: Kant, Enlightenment, and Courage}\par\vspace{0.45em}
\rule{0.82\textwidth}{0.4pt}\par\vspace{0.7em}
\end{center}

\section*{Vocabulary Preview}
\begin{multicols}{2}
\begin{itemize}[leftmargin=*, itemsep=0.18em]
\item \textbf{Enlightenment}: The process of coming out of intellectual dependence and using one's own reason.
\item \textbf{Immaturity}: For Kant, the inability or refusal to use one's understanding without another person's direction.
\item \textbf{Self-incurred}: Brought upon oneself; not forced entirely from outside.
\item \textbf{Tutelage}: Being under another person's guidance or control, like a pupil or dependent child.
\item \textbf{Resolution}: Firm decision, determination, or settled purpose.
\item \textbf{Courage}: The willingness to face difficulty or risk for the sake of what is right.
\item \textbf{Guardian}: A person or institution that thinks, decides, or judges for others.
\item \textbf{Public use of reason}: Speaking or writing as a thinker before the public, offering reasons and criticism.
\item \textbf{Private use of reason}: Reasoning within an assigned role or office where obedience may be required.
\item \textbf{Prejudice}: A settled judgment accepted before careful thought.
\item \textbf{Dignity}: The worth proper to a rational human being, not a mere tool or machine.
\end{itemize}
\end{multicols}

\section*{Introduction}
Immanuel Kant lived from 1724 to 1804 in Prussia, in what is now Germany. In 1784 he answered a public question: \textit{What is enlightenment?} Kant's answer belongs in this unit because he does not treat shadows as only a mistake of the eyes, a beautiful temptation, a distorted mind, or a disordered love. He asks whether people remain in shadows because maturity is difficult. It is often easier to let someone else think, judge, and decide.

For Kant, the problem is not that human beings lack reason. The problem is that many lack the courage and discipline to use reason without being led. As you read, ask: \textbf{Do people remain in shadows because truth requires the courage to think for oneself?}

\focusbox{\textbf{Source note:} Adapted and modernized from Immanuel Kant, ``An Answer to the Question: What Is Enlightenment?'' (1784), using the public-domain English text preserved by the Internet Modern History Sourcebook and LibreTexts. This focused version preserves Kant's definition of enlightenment, his warning about guardians, his public/private use of reason distinction, and his claim that enlightenment requires freedom.}

\begin{multicols}{2}
\raggedcolumns
\section*{Reading}

\subsection*{1. What Enlightenment Is}
Enlightenment is the human being's release from self-incurred immaturity. Immaturity is the inability to use one's own understanding without direction from another. This immaturity is self-incurred when its cause is not a lack of understanding, but a lack of resolution and courage to use that understanding without another person's direction.

\textit{Sapere aude}: dare to know. Have courage to use your own reason. That is the motto of enlightenment.

Laziness and cowardice are the reasons why so many people, even after nature has long set them free from outside direction, still remain immature for life. These are also the reasons why it is so easy for others to set themselves up as guardians over them.

It is so comfortable not to be mature. If I have a book that understands for me, a pastor who has a conscience for me, a doctor who decides my diet for me, and so on, I do not need to trouble myself. I do not need to think, if only I can pay. Others will easily take over that tiresome work for me.

\subsection*{2. Why Guardians Want People Afraid}
The step toward maturity is held to be very dangerous by most people. The guardians who have so kindly taken charge of them make sure of this. First, the guardians make their tame animals unable to walk freely. They make sure these quiet creatures dare not take a single step without the cart-harness to which they are tied. Then the guardians show them the danger that threatens if they try to walk alone.

But the danger is not really so great. By falling a few times, people would finally learn to walk. Yet one example of failure makes them timid and usually frightens them away from further attempts.

For one individual to work himself out of immaturity is very difficult, because that condition has almost become his nature. He has even grown fond of it. For the present, he is truly unable to use his own reason, because no one has ever allowed him to try.

Rules, formulas, and ready-made systems are mechanical tools for using reason---or rather, for misusing natural gifts. They become the chains of permanent immaturity. Whoever throws them off makes only an uncertain leap over a narrow ditch, because he is not used to free movement. Therefore only a few have succeeded, by exercising their own minds, in freeing themselves from immaturity and walking steadily.

\subsection*{3. Why a Public Can Become Enlightened}
It is more possible for a whole public to enlighten itself. Indeed, if freedom is granted, enlightenment is almost sure to follow. There will always be some independent thinkers, even among the established guardians of the great masses. After they throw off the yoke of immaturity from their own shoulders, they will spread the spirit of rational appreciation for human worth and for every person's calling to think for himself.

Still, the public can reach enlightenment only slowly. A revolution may overthrow personal despotism or greedy and tyrannical oppression, but it will never by itself produce a true reform in ways of thinking. New prejudices will serve just as well as old ones to harness the great unthinking masses.

For this enlightenment, nothing is required except freedom---and the most harmless kind of freedom that can be called by that name: the freedom to make public use of one's reason in all matters.

But from all sides I hear the command, ``Do not argue!'' The officer says, ``Do not argue, but drill!'' The tax collector says, ``Do not argue, but pay!'' The cleric says, ``Do not argue, but believe!'' Everywhere there is restriction on freedom.

\subsection*{4. Public and Private Use of Reason}
Which restriction is an obstacle to enlightenment, and which is not? I answer: the public use of one's reason must always be free, and it alone can bring about enlightenment among human beings. The private use of reason, however, may often be narrowly restricted without greatly hindering the progress of enlightenment.

By public use of reason, I mean the use a person makes of reason as a scholar before the reading public. By private use of reason, I mean the use a person makes of reason within a particular civil post or office entrusted to him.

Many affairs of the community require a certain order. In those affairs, some members must act as part of a mechanism so that public goals can be accomplished. There, argument is often not allowed; one must obey. But the same person, considered as a member of the whole community, or even of a society of world citizens, may write and speak to the public as a scholar. In that role he may argue without damaging the work assigned to him.

For example, it would be ruinous for an officer on duty to debate a command given by a superior. He must obey. But as a scholar, that same officer should not be denied the right to publish remarks about errors in military service and to submit them to public judgment. A citizen cannot refuse to pay taxes that have been imposed. Yet the same citizen does not act against duty if, as a scholar, he publicly expresses thoughts about the injustice or poor design of those taxes.

So also a clergyman must teach his congregation according to the creed of the church he serves. In that office, he acts under an assigned trust. But as a scholar, he has full freedom---and even the calling---to communicate carefully tested thoughts to the public about faults in that creed and proposals for improving religious and church life.

\subsection*{5. Can an Age Bind the Future?}
A group of church leaders might agree to bind themselves to one unchangeable creed, so that each member would guard the others and the whole institution would hold the people under a fixed doctrine. But such a contract would be absolutely null and void if it were meant to shut off all further enlightenment from the human race.

One age cannot bind itself and arrange conditions so that a later age cannot expand its knowledge, purify itself from errors, and make progress in enlightenment. That would be a crime against human nature, whose original vocation lies precisely in such progress. Later generations are entirely justified in rejecting such decisions as unauthorized and wicked.

\subsection*{6. An Age of Enlightenment}
If we are asked, ``Do we now live in an enlightened age?'' the answer is no. But we do live in an age of enlightenment. As things now stand, much is still lacking before human beings as a whole can use their own understanding well and confidently in religious matters without outside direction. Yet there are clear signs that the field has been opened for people to work freely in that direction, and that the obstacles to general enlightenment are gradually becoming fewer.

I have placed the main point of enlightenment---the human escape from self-incurred immaturity---chiefly in matters of religion, because rulers have little interest in playing guardian over the arts and sciences, and because religious immaturity is not only harmful but deeply degrading.

Only a ruler who is himself enlightened, who is not afraid of shadows, and who can preserve public peace may say: ``Argue as much as you will, and about what you will, only obey.'' In this strange course of human affairs, a greater degree of civil freedom may give less room for freedom of mind, while a lower degree of civil freedom may give the mind room to expand. But when nature has uncovered the seed she most cares for---the human tendency and calling to free thought---that free thought gradually works back upon the character of the people. At last it even affects the principles of government, which finds it useful to treat human beings according to their dignity, not as machines.

\section*{Reading Focus}
\begin{enumerate}[leftmargin=*, itemsep=0.2em]
\item Underline Kant's definition of enlightenment.
\item Circle one sentence showing why people choose dependence or immaturity.
\item Put a G in the margin beside one sentence about guardians and fear.
\item Put a P beside Kant's explanation of the public use of reason.
\item Box one sentence that best answers this question: Why might people remain in shadows even when they have the ability to think?
\end{enumerate}

\section*{Obstacle Map}
Draw a simple obstacle map showing what keeps people from using their own reason. Include at least four obstacles from Kant: laziness, cowardice, guardians, fear of failure, rules/formulas, habit, or lack of freedom. For each obstacle, write one sentence explaining how it keeps someone in shadows.

\section*{Enlightenment Claim}
Answer in one clear paragraph: \textbf{What makes it hard to use your own reason?} State a claim, use one specific moment from Kant, and explain how his answer adds to Plato, Homer, Bacon, or Augustine.

\section*{Handout Summary}
Finish by writing a short summary of this handout. In three to five sentences, explain Kant's definition of enlightenment, why people remain immature, and why freedom to use reason publicly matters.

\end{multicols}
\end{document}
